\section{和议年代的财政国家}

1163年符离受挫、1164年隆兴和议之后，南宋进入一段常被称为“和平”的时期；1163--1189年的连续在位纪年提醒我们，和平首先是一套要不断支付的制度。两淮、四川和江汉的防务，临安与地方之间的漕运，军费、盐茶与纸币管理，都使边界并未退出日常行政。和约的条款或会要中的财政规定，不能自动证明一项征收在所有地区、所有户等中同样实现。

1150年前后会子管理的调整，以及1161年采石等战事，已表明货币、军事与江南供给彼此相连。研究会子时，发行额不等于流通量，更不等于物价或家庭债务；必须辨认币制、地域、期限和记录机构。研究军费时，同样要分开中央预算、地方催科、实物运输和边军实际所得。只有把这些层次拆开，才不会用“富庶江南”抹平不同地区的压力。

\section{地方材料所见的国家边缘}

南宋契约、账籍、讼牍、碑刻、寺院材料和港口考古可以让土地、租佃、借贷、婚姻、差役与工艺进入叙事，但保存最密集的江南、福建、两浙材料不能代替四川、两广和边区。文书先要辨原件还是抄件、形成地点、当事人身份与用途；一纸买卖契约证明的是一次交易，不是普遍土地制度。

1173年前后的书院、讲学和地方活动也属于这种制度社会的一部分。它们连接士人、地方行政和文本传播，却不能由一位名儒或一座书院推及全体乡村。南宋能在1160年代以后维持政权，正因为地方征收、市场、交通、军事与知识网络在不同尺度上继续运作；同样也因为它们并不完全服从临安的设计。
